\documentclass[10pt,a4paper]{article}
\usepackage[utf8]{inputenc}
\usepackage[T1]{fontenc}
\usepackage[a4paper,landscape,margin=1.35cm]{geometry}
\usepackage{array}
\usepackage{longtable}

\setlength{\parindent}{0pt}
\setlength{\parskip}{0.45em}
\setlength{\tabcolsep}{3.5pt}
\renewcommand{\arraystretch}{1.18}

\newcolumntype{L}[1]{>{\raggedright\arraybackslash}p{#1}}
\newcommand{\code}[1]{\texttt{\detokenize{#1}}}

\begin{document}

{\Large \textbf{Auditoria dos datasets ERA5 em \code{Datasets/nc_files/}}}

\textbf{Data da auditoria:} 13/07/2026.

\textbf{Escopo:} leitura de metadados dos arquivos NetCDF presentes em
\code{Datasets/nc_files/}. A auditoria usa cabeçalhos, dimensões, coordenadas e
atributos das variáveis; os campos meteorológicos completos não foram
materializados em memória. A pasta \code{Datasets/nc_files/} é tratada no projeto
como localização legada de NetCDF, portanto este documento registra o acervo sem
alterar os datasets.

\section*{Resumo dos arquivos}

\footnotesize
\begin{longtable}{|L{4.3cm}|L{2.2cm}|L{4.4cm}|L{5.9cm}|L{5.3cm}|L{1.6cm}|}
\hline
\textbf{Arquivo} &
\textbf{Variável principal} &
\textbf{Descrição e unidade} &
\textbf{Dimensões principais} &
\textbf{Data inicial e final} &
\textbf{Tamanho} \\
\hline
\endfirsthead
\hline
\textbf{Arquivo} &
\textbf{Variável principal} &
\textbf{Descrição e unidade} &
\textbf{Dimensões principais} &
\textbf{Data inicial e final} &
\textbf{Tamanho} \\
\hline
\endhead

\code{cape_88-26.nc} &
\code{cape} &
Convective available potential energy; \code{J kg**-1} &
3 dimensões: \code{time}=335059, \code{latitude}=41, \code{longitude}=61 &
1988-01-01 00:00 até 2026-03-22 18:00; amostra temporal de 1 h &
3.36 GB \\
\hline

\code{era5_precipitation_80-26.nc} &
\code{tp} &
Total precipitation; \code{m} &
4 dimensões: \code{time}=33653, \code{step}=12, \code{latitude}=41, \code{longitude}=61 &
1979-12-31 18:00 até 2026-01-24 18:00; amostra em \code{time} de 12 h, com 12 passos em \code{step} &
4.04 GB \\
\hline

\code{era5_sh_00-25.nc} &
\code{q} &
Specific humidity; \code{kg kg**-1} &
4 dimensões: \code{time}=227928, \code{isobaricInhPa}=2, \code{latitude}=41, \code{longitude}=61 &
2000-01-01 00:00 até 2025-12-31 23:00; amostra temporal de 1 h &
4.56 GB \\
\hline

\code{era5_wind_14-25.nc} &
\code{u}, \code{v} &
U/V component of wind; \code{m s**-1} &
4 dimensões: \code{time}=105192, \code{isobaricInhPa}=2, \code{latitude}=41, \code{longitude}=61 &
2014-01-01 00:00 até 2025-12-31 23:00; amostra temporal de 1 h &
4.21 GB \\
\hline

\code{temp_99-25.nc} &
\code{d2m}, \code{t2m} &
2 metre dewpoint temperature e 2 metre temperature; \code{K} &
3 dimensões: \code{time}=236688, \code{latitude}=41, \code{longitude}=61 &
1999-01-01 00:00 até 2025-12-31 23:00; amostra temporal de 1 h &
4.74 GB \\
\hline

\code{vv_94-25.nc} &
\code{w} &
Vertical velocity; \code{Pa s**-1} &
4 dimensões: \code{time}=280512, \code{isobaricInhPa}=3, \code{latitude}=41, \code{longitude}=61 &
1994-01-01 00:00 até 2025-12-31 23:00; amostra temporal de 1 h &
8.42 GB \\
\hline

\end{longtable}
\normalsize

\section*{Variáveis de dados}

\footnotesize
\begin{longtable}{|L{4.3cm}|L{1.3cm}|L{4.7cm}|L{2.0cm}|L{5.5cm}|L{1.2cm}|L{1.4cm}|}
\hline
\textbf{Arquivo} &
\textbf{Var.} &
\textbf{Nome longo} &
\textbf{Unidade} &
\textbf{Standard name} &
\textbf{Dims} &
\textbf{Tipo} \\
\hline
\endfirsthead
\hline
\textbf{Arquivo} &
\textbf{Var.} &
\textbf{Nome longo} &
\textbf{Unidade} &
\textbf{Standard name} &
\textbf{Dims} &
\textbf{Tipo} \\
\hline
\endhead

\code{cape_88-26.nc} & \code{cape} & Convective available potential energy & \code{J kg**-1} & \code{unknown} & 3 & \code{float32} \\
\hline
\code{era5_precipitation_80-26.nc} & \code{tp} & Total precipitation & \code{m} & \code{unknown} & 4 & \code{float32} \\
\hline
\code{era5_sh_00-25.nc} & \code{q} & Specific humidity & \code{kg kg**-1} & \code{specific_humidity} & 4 & \code{float32} \\
\hline
\code{era5_wind_14-25.nc} & \code{u} & U component of wind & \code{m s**-1} & \code{eastward_wind} & 4 & \code{float32} \\
\hline
\code{era5_wind_14-25.nc} & \code{v} & V component of wind & \code{m s**-1} & \code{northward_wind} & 4 & \code{float32} \\
\hline
\code{temp_99-25.nc} & \code{d2m} & 2 metre dewpoint temperature & \code{K} & \code{unknown} & 3 & \code{float32} \\
\hline
\code{temp_99-25.nc} & \code{t2m} & 2 metre temperature & \code{K} & \code{unknown} & 3 & \code{float32} \\
\hline
\code{vv_94-25.nc} & \code{w} & Vertical velocity & \code{Pa s**-1} & \code{lagrangian_tendency_of_air_pressure} & 4 & \code{float32} \\
\hline

\end{longtable}
\normalsize

\section*{Grade espacial e coordenadas}

\footnotesize
\begin{longtable}{|L{4.3cm}|L{4.5cm}|L{4.5cm}|L{4.6cm}|L{5.3cm}|}
\hline
\textbf{Arquivo} &
\textbf{Latitude} &
\textbf{Longitude} &
\textbf{Coordenada vertical ou superfície} &
\textbf{Coordenadas auxiliares presentes} \\
\hline
\endfirsthead
\hline
\textbf{Arquivo} &
\textbf{Latitude} &
\textbf{Longitude} &
\textbf{Coordenada vertical ou superfície} &
\textbf{Coordenadas auxiliares presentes} \\
\hline
\endhead

\code{cape_88-26.nc} &
-35.0 a -25.0; 41 pontos; passo -0.25 grau &
-60.0 a -45.0; 61 pontos; passo 0.25 grau &
\code{surface} &
\code{number}, \code{step}, \code{valid_time} \\
\hline

\code{era5_precipitation_80-26.nc} &
-35.0 a -25.0; 41 pontos; passo -0.25 grau &
-60.0 a -45.0; 61 pontos; passo 0.25 grau &
\code{surface}; dimensão \code{step}=12 &
\code{number}, \code{step}, \code{valid_time} \\
\hline

\code{era5_sh_00-25.nc} &
-35.0 a -25.0; 41 pontos; passo -0.25 grau &
-60.0 a -45.0; 61 pontos; passo 0.25 grau &
\code{isobaricInhPa}=2 níveis: 1000, 850 hPa &
\code{number}, \code{step}, \code{valid_time} \\
\hline

\code{era5_wind_14-25.nc} &
-35.0 a -25.0; 41 pontos; passo -0.25 grau &
-60.0 a -45.0; 61 pontos; passo 0.25 grau &
\code{isobaricInhPa}=2 níveis: 850, 500 hPa &
\code{number}, \code{step}, \code{valid_time} \\
\hline

\code{temp_99-25.nc} &
-35.0 a -25.0; 41 pontos; passo -0.25 grau &
-60.0 a -45.0; 61 pontos; passo 0.25 grau &
\code{surface} &
\code{number}, \code{step}, \code{valid_time} \\
\hline

\code{vv_94-25.nc} &
-35.0 a -25.0; 41 pontos; passo -0.25 grau &
-60.0 a -45.0; 61 pontos; passo 0.25 grau &
\code{isobaricInhPa}=3 níveis: 850, 500, 300 hPa &
\code{number}, \code{step}, \code{valid_time} \\
\hline

\end{longtable}
\normalsize

\section*{Observações principais}

\begin{itemize}
    \item Todos os arquivos declaram convenção \code{CF-1.7} e instituição
    \code{European Centre for Medium-Range Weather Forecasts}.
    \item A grade horizontal é comum a todos os arquivos: 41 latitudes por 61
    longitudes, cobrindo aproximadamente -35 a -25 graus de latitude e -60 a
    -45 graus de longitude, com resolução nominal de 0.25 grau.
    \item A maioria dos arquivos tem amostragem horária na coordenada
    \code{time}. A exceção é \code{era5_precipitation_80-26.nc}, que combina
    \code{time} a cada 12 horas com \code{step}=12; para esse caso, a coordenada
    derivada \code{valid_time} deve ser considerada ao alinhar precipitação com
    variáveis horárias.
    \item As variáveis em níveis de pressão usam \code{isobaricInhPa}: umidade
    específica tem níveis 1000 e 850 hPa; vento tem níveis 850 e 500 hPa; e
    velocidade vertical tem níveis 850, 500 e 300 hPa.
    \item Os arquivos \code{temp_99-25.nc} e \code{era5_wind_14-25.nc} carregam
    duas variáveis meteorológicas cada; rotinas de extração devem tratar ambas
    explicitamente para não perder canais de entrada.
\end{itemize}

\end{document}
